% report-template.tex — kerangka laporan pengujian keamanan
% Compile: pdflatex report-template  (jalankan 3x agar ToC dan referensi stabil)
% Catatan: template ini kerangka minimal. Ganti setiap <...> dan hapus komentar.
\documentclass[11pt,a4paper]{article}

\usepackage[T1]{fontenc}
\usepackage[utf8]{inputenc}
\usepackage[a4paper,margin=2.2cm]{geometry}
\usepackage{graphicx}
\usepackage{longtable}
\usepackage{booktabs}
\usepackage{xcolor}
\usepackage{listings}
\usepackage{enumitem}
\usepackage{fancyhdr}
\usepackage{titlesec}
\usepackage{caption}
\usepackage{setspace}
\usepackage[hidelinks]{hyperref}

\graphicspath{{notes/}{assets/}}

% --- Badge severity -------------------------------------------------
\definecolor{sevcritical}{HTML}{8B0000}
\definecolor{sevhigh}{HTML}{D9534F}
\definecolor{sevmed}{HTML}{F0AD4E}
\definecolor{sevlow}{HTML}{5BC0DE}
\definecolor{sevinfo}{HTML}{777777}

\newcommand{\severity}[2]{\colorbox{#1}{\color{white}\textbf{\small #2}}}
\newcommand{\riskcritical}[1]{\severity{sevcritical}{#1}}
\newcommand{\riskhigh}[1]{\severity{sevhigh}{#1}}
\newcommand{\riskmed}[1]{\severity{sevmed}{#1}}
\newcommand{\risklow}[1]{\severity{sevlow}{#1}}
\newcommand{\riskinfo}[1]{\severity{sevinfo}{#1}}

% --- Listing --------------------------------------------------------
\lstset{
  basicstyle=\ttfamily\small,
  breaklines=true,
  frame=single,
  columns=fullflexible,
  keepspaces=true,
  showstringspaces=false
}

\pagestyle{fancy}
\fancyhf{}
\fancyhead[L]{Laporan Pengujian Keamanan}
\fancyfoot[C]{\thepage}

\title{Laporan Pengujian Keamanan\\\large <Aplikasi / Target>}
\author{<Tim>}
\date{<dd MMMM yyyy>}

\begin{document}
\maketitle
\tableofcontents
\newpage

\section{Ringkasan Pengujian}
% Paragraf non-teknis: nomor tiket, tujuan, target, waktu, reverifikasi.

\section{Ringkasan Temuan Keamanan}
\begin{table}[h]
\centering
\begin{tabular}{lr}
\toprule
\textbf{Severity} & \textbf{Jumlah} \\
\midrule
CRITICAL & 0 \\
HIGH     & 0 \\
MEDIUM   & 0 \\
LOW      & 0 \\
INFO     & 0 \\
\midrule
\textbf{Total} & \textbf{0} \\
\bottomrule
\end{tabular}
\end{table}

\section{Ringkasan Teknis}
\subsection{Waktu Pengujian}
\subsection{Cakupan Poin Uji}
% Sebutkan juga yang TIDAK diuji — lihat bagian Gotchas di SKILL.md.
\subsection{Metode Pengujian}
\subsection{Temuan Utama}
\subsection{Rekomendasi Utama}

\section{Detail Pengujian}
% Satu subsection per temuan, urut dari severity tertinggi.
\subsection{<Judul temuan>}

\begin{longtable}{@{}p{3.4cm}p{11cm}@{}}
\toprule
\textbf{Severity}   & \riskhigh{HIGH} \\
\textbf{CVSS v3.1}  & 8.8 (\texttt{AV:N/AC:L/PR:L/UI:N/S:U/C:H/I:H/A:H}) \\
\textbf{CWE}        & CWE-639 \\
\textbf{Komponen}   & \texttt{/api/v1/users/\{id\}} \\
\bottomrule
\end{longtable}

\paragraph{Deskripsi.} <Apa yang terjadi, dalam dua sampai empat kalimat.>

\paragraph{Dampak bisnis.} <Konsekuensi nyata bagi organisasi.>

\paragraph{Bukti (reproducible).}
\begin{lstlisting}
curl -s -H "Authorization: Bearer $TOKEN" \
  https://target.tld/api/v1/users/1042
\end{lstlisting}

\paragraph{Rekomendasi.} <Perbaikan spesifik, bukan "tingkatkan keamanan".>

\section{Kesimpulan}

\appendix
\section{Lampiran}
\subsection{Role dan Scope Mapping}
\subsection{API Endpoint Mapping}
\subsection{Daftar Hash Bukti}

\end{document}
